

The oriented manifold
$M_{\ell,m}=\ell\mathbb{HP}^2\#m\overline{\mathbb{HP}^2}$
admits an orientation-compatible almost complex structure exactly
when $\ell=2m+1$ and $m$ is odd.  In particular $M_{3,1}$ is
almost complex and has the invariants \eqref{maps:eq:F.2}.

\begin{proof}
This is the exact criterion of Mandal's Corollary 1.4
\cite{ffMandal}; we read both that statement and its proof in
Section 5.  Here is its application in the normalization needed
here.  The tangent invariants in Mandal's normalized basis are
$p_1=2\sum x_i$, $p_2=7(\ell-m)u$ and
$e=(\ell+m+2)u$.  The stable-complex congruence of his
Theorem 1.1 becomes
$7(\ell-m)\equiv\ell-m\pmod4$, or $\ell-m$ even.
The rank-eight criterion of his Theorem 1.3(iii) becomes
$6(\ell-m)=2(\ell+m+2)$, or $\ell=2m+1$.
Together they say precisely that $m$ is odd.  At $(3,1)$ we
obtain $\chi=6$, $\sigma=2$, $p_1^2=8$ and $c_2^2=2$.
An integrable complex structure is not asserted by that theorem.
The later preprint arXiv:2610.02147 appeared in the literature
search, but its full text could not be opened in this run and
is not used as an authority.
\end{proof}


